\chapter{Applicazioni client-side con React}

\section{Architetture web: i flussi di informazione}

Elementi principali di un'architettura web: \textbf{Data Persistence},
\textbf{Application Business Logic}, \textbf{User Interface}, pi\`u eventuali
\textbf{Third Party Services}. I flussi tipici:
\begin{itemize}
  \item \textbf{Flow 1 -- retrieval + presentation}: dai dati persistenti, attraverso la BL,
        alla UI (la BL elabora i dati ma pu\`o elaborare anche elementi di
        \emph{presentazione});
  \item \textbf{Flow 2a -- user action, UI-only}: l'azione tocca solo la UI (spostare un
        elemento, selezionare);
  \item \textbf{Flow 2b -- user action, non-persistent}: modifica ai dati non ancora da
        salvare, o azione che influenza la presentazione ma non i dati;
  \item \textbf{Flow 2c -- user action con cambiamento persistente}: l'azione arriva fino
        alla persistenza;
  \item \textbf{Flow 3 -- external update}: come il flow 1, ma originato da un \emph{evento
        esterno} ad applicazione gi\`a attiva, non dall'utente.
\end{itemize}
La logica pu\`o essere distribuita: \textbf{client-side logic} (BL sul client),
\textbf{server-side logic} (BL sul server), \textbf{two-sided logic} (divisa fra i due lati).

\section{React: caratteristiche principali}

\begin{definizione}[React]
\textbf{React} \`e propriamente una \emph{libreria} (non un framework): pu\`o essere usato da
solo o dentro un framework che lo supporta (``il'' framework di React \`e Next.JS). Nel corso
si usa come libreria in abbinamento al tool \textbf{Vite} che, con npm, facilita
configurazione e gestione del progetto. React supporta lo sviluppo di front-end con
\textbf{UI orientate ai componenti} (come i competitor Angular e Vue, che per\`o sono
framework a tutto tondo) e di \textbf{Single Page Application}: applicazioni client-side
costituite da un'\emph{unica pagina HTML} il cui DOM viene costantemente aggiornato
dall'applicazione stessa. A differenza di altri framework, React \textbf{non \`e
object-oriented}: \`e fortemente basato sugli aspetti \emph{funzionali} di JS/TS.
\end{definizione}

\textbf{Libreria vs framework.} Una \emph{libreria} \`e un insieme di package importabili di
cui si usano le funzionalit\`a; React \`e complessa (include un'estensione del linguaggio che
richiede compilazione, e un \emph{runtime} di oggetti e funzioni che gestiscono i
componenti), ma sono sempre funzionalit\`a usate \emph{dal nostro codice}. Un vero
\emph{framework} realizza qualcosa di simile a una ``virtual machine'': un ambiente di
esecuzione dentro l'ambiente di esecuzione, che carica ed esegue il nostro codice in un
contesto altamente modificato, compiendo operazioni ``dietro le quinte'' (es.: Angular
monitora tutti i cambiamenti delle propriet\`a visualizzate; Vue crea copie ``addizionate''
degli oggetti). I framework sono pi\`u completi e coprono anche aspetti oltre la UI; in React
questi aspetti sono demandati a librerie separate.

\section{Componenti UI e gerarchie}

\begin{definizione}[Componente UI]
Un componente UI \`e caratterizzato da:
\begin{itemize}
  \item un \textbf{aspetto grafico} (il \emph{template}): come viene visualizzato -- sul web
        client-side, una combinazione di HTML e CSS;
  \item una \textbf{parte del Data Model} visualizzata e/o modificata attraverso il
        componente (esistono anche componenti che mostrano informazioni statiche o servono
        solo alla navigazione/interazione);
  \item un \textbf{comportamento}: quali operazioni pu\`o compiere l'utente attraverso il
        componente e come si svolge l'interazione.
\end{itemize}
\end{definizione}

\begin{importante}[Genitori e figli: sono le istanze a essere figlie]
Un componente pu\`o essere costruito a partire da (istanze di) altri componenti: \`e detto
\textbf{genitore} di quelle istanze, che sono i suoi \textbf{figli}. Esempio classico: un
menu di navigazione, in cui la singola voce \`e un componente istanziato pi\`u volte (stesso
aspetto e comportamento, contenuto diverso), e il menu costruisce altro sul comportamento
delle voci (es. garantire che una sola sia selezionata). \textbf{Non \`e il componente in
astratto a essere figlio, ma le sue istanze}: istanze diverse dello stesso componente possono
avere genitori diversi. La definizione di un componente non include ``chi lo conterr\`a'', ma
solo come chi lo conterr\`a potr\`a comunicare con esso; un componente in astratto pu\`o
invece dirsi \emph{genitore}, perch\'e la sua definizione specifica come \`e costruito.
\end{importante}

\subsection{Componenti dinamici: show/hide vs mount/unmount}
I componenti cambiano aspetto in base a dati e interazioni -- incluso il ``set'' di figli
mostrati (un menu espandibile: compresso non mostra sotto-voci). La \emph{gerarchia} di cui
parliamo \`e quella \textbf{statica} (tutti i figli dichiarati, indipendentemente da quelli
mostrati). Due modi per far scomparire/riapparire un figlio:
\begin{itemize}
  \item \textbf{show / hide}: cambiare la classe CSS aggiungendo/togliendo
        \texttt{display:none}; il figlio resta nell'HTML del genitore, solo non visibile.
        Preferito quando il figlio va acceso/spento \emph{di frequente e rapidamente};
  \item \textbf{mount / unmount}: aggiungerlo/toglierlo dal DOM. Preferito per
        sotto-componenti \emph{sostanziosi} mostrati solo occasionalmente, o quando la
        struttura dei figli \`e generata dai dati e i dati cambiano.
\end{itemize}

\subsection{Single Page Application}
Le diverse ``schermate'' dell'app sono tutte figlie di un macro-componente principale,
\textbf{App}, che monta \emph{un solo figlio alla volta}: quando l'utente ``naviga'', App
smonta un figlio e ne monta un altro -- \emph{non} c'\`e mai una vera navigazione (che
provocherebbe un reload e il reset dell'applicazione); la pagina non contiene mai veri link
ipertestuali. Due modi di realizzarlo:
\begin{itemize}
  \item \textbf{gestione manuale}: menu come insiemi di ``pulsanti'' che montano/smontano i
        sotto-componenti; semplice ma macchinosa su app articolate;
  \item \textbf{router}: una libreria che permette di definire \emph{route} (analoghe a
        quelle server-side) con i componenti da montare per ciascuna; il router
        \emph{intercetta} la navigazione e configura la UI. Vantaggio: preserva all'utente
        l'illusione di navigare fra pagine (la barra del browser mostra la route attiva
        anzich\'e la URL dell'unica vera pagina).
\end{itemize}

\section{Il pattern MVVM}

\begin{esamebox}
\emph{``Illustrare il pattern architetturale MVVM per come \`e stato presentato nel contesto
di React.''}
\end{esamebox}

\begin{definizione}[Model--View--ViewModel]
\begin{itemize}
  \item \textbf{Model} (modello dati): i dati che l'app elabora e i vincoli che ne regolano
        fruizione ed elaborazione;
  \item \textbf{View}: ci\`o che l'utente vede e con cui interagisce (\emph{il DOM});
  \item \textbf{View-Model}: il \emph{comportamento dei componenti} -- legge i dati e li
        presenta nella view; riceve gli eventi UI, richiede modifiche ai dati, rilegge i dati
        aggiornati e aggiorna la presentazione.
\end{itemize}
Il nucleo del pattern \`e il \textbf{legame (binding)} fra View-Model e View, che nei
framework \`e \emph{parzialmente automatizzato}, in due direzioni:
\begin{itemize}
  \item \textbf{property binding}: la View riflette i dati del View-Model; ogni framework ha
        la sua strategia per accorgersi di quando qualcosa cambia nel view-model e aggiornare
        la view;
  \item \textbf{event binding}: dalla View verso il View-Model; basato sul principio di
        callback ed event listener, associa una funzione a un evento UI.
\end{itemize}
\end{definizione}

\section{I componenti React: il ViewModel del MVVM}

\begin{definizione}[Componente React]
Un componente in React \`e una \textbf{funzione} (detta \emph{funzione di rendering}) che
costruisce il sottoalbero del DOM che costituir\`a la \emph{view} del componente.
\begin{itemize}
  \item I componenti sono scritti in \textbf{JSX/TSX}, che estende JS/TS permettendo di
        \emph{interpolare} HTML e codice (e di importare direttamente i CSS).
  \item La funzione genera la view \emph{a partire dai dati rilevanti} per il componente;
        il componente \emph{dichiara} da quali dati dipende la sua view, e ogni volta che
        quei dati cambiano React invoca di nuovo la funzione, rigenerando la view e
        sostituendo l'elemento nel DOM: \`e la base del \textbf{property binding} in React.
  \item Nel generare la view, la funzione specifica anche le \emph{callback} che rispondono
        agli eventi del DOM, stabilendo il comportamento; le callback possono modificare i
        dati del componente (provocandone il re-rendering): \`e la base
        dell'\textbf{event binding}.
  \item A ogni componente React associa un \textbf{tag} con lo stesso nome, usabile nella
        view di altri componenti come un tag HTML: cos\`i si definisce la \emph{gerarchia}.
\end{itemize}
\end{definizione}

Il nome della funzione diventa il nome del componente (esportabile per usarlo altrove);
l'output deve essere un oggetto \textbf{ReactElement} (da un'espressione JSX/TSX); il tipo
dell'oggetto-funzione \`e \texttt{FunctionComponent}.
\begin{attenzione}[I componenti devono essere funzioni pure]
React richiede che l'output della funzione di rendering dipenda \emph{solo dall'input}, e che
la funzione \emph{non alteri} strutture dati fuori dal proprio scope. Questo vale per la
funzione di rendering in s\'e, \textbf{non} per le callback usate al suo interno!
\end{attenzione}

\subsection{Espressioni JSX/TSX}
\begin{lstlisting}[language=JavaScript]
const elem = <p>Pagina creata da {utente.nome} {utente.cognome}</p>
\end{lstlisting}
L'HTML \`e riportato direttamente nel codice (senza virgolette); fra graffe
(\textbf{interpolazione}) ricompare JS/TS: il codice interpolato deve risolversi in
un'espressione, il cui valore completa l'HTML in quel punto.
\begin{itemize}
  \item Un'espressione JSX/TSX deve risultare in \textbf{esattamente un sottoalbero} del DOM
        (una sola radice): \texttt{<h1>...</h1><p>...</p>} non \`e valida. Se servono
        pi\`u radici, si racchiudono nel \textbf{tag vuoto} \texttt{<>\dots</>}: si ottiene
        un ReactElement valido, e il tag vuoto viene tolto all'inserimento nel DOM.
  \item Si usa il normale HTML con accorgimenti: \emph{tutti i tag vanno chiusi} (anche
        quelli senza contenuto: \texttt{<img src="..."/>}); gli attributi in kebab-case
        diventano \emph{camelCase} (\texttt{data-value} $\rightarrow$ \texttt{dataValue}:
        diventano propriet\`a di oggetti JS, e i nomi di propriet\`a non ammettono il
        trattino); \texttt{class} diventa \textbf{className} (\texttt{class} \`e keyword
        riservata in JS/TS).
\end{itemize}

\textbf{Dove pu\`o comparire l'interpolazione:}
\begin{enumerate}
  \item come intero contenuto di un tag: \texttt{<ul>\{espressione\}</ul>};
  \item come porzione di testo: \texttt{<p>\{nome\} ha vinto \{x+y\} euro!</p>};
  \item come valore di un attributo: \texttt{<div className=\{val\}>}.
\end{enumerate}
Nei primi due casi il valore pu\`o anche essere un ReactElement \emph{o un array di
ReactElement}, che vengono inseriti in quel punto del DOM:
\begin{lstlisting}[language=JavaScript]
const lista = [<li>Uno</li>,<li>Due</li>,<li>Tre</li>];
return <ul>{lista}</ul>;
\end{lstlisting}
Il terzo caso \`e quello che permette di associare comportamento agli eventi UI (ci\`o che
senza React si faceva con i listener): per ogni evento \`e definito un attributo
corrispondente (\texttt{onClick}, \texttt{onMouseDown}, \texttt{onKeyPress}\dots) a cui si
assegna il listener per interpolazione:
\begin{lstlisting}[language=JavaScript]
<button onClick={(e) => myFunction(e, ...)}>PREMI!</button>
\end{lstlisting}

\section{Le props: l'input di un componente}

\begin{definizione}[Props]
Un componente pu\`o dichiarare dei parametri che il genitore, nell'istanziarlo, dovr\`a
valorizzare: collettivamente sono le \textbf{props}. Il componente prende come argomento un
oggetto \texttt{props}; la sua ``forma'' va dichiarata in una \emph{interface} che ne
costituisce il tipo. Il genitore assegna i valori al momento dell'istanziazione: i nomi delle
props diventano \emph{attributi} del tag del componente.
\begin{lstlisting}[language=JavaScript]
interface BookProps { bIS: BookInBookshelf; }
function Book({bIS}: BookProps): ReactElement {...}

function Bookshelf(...): ReactElement {
  const bookElems = books.map(book => <Book bIS={book} />);
  return (<div className="libreria">{bookElems}</div>);
}
\end{lstlisting}
\end{definizione}

Vantaggi evidenti: \textbf{riuso} (stesso componente su dati diversi),
\textbf{configurabilit\`a} (il genitore determina aspetti della view del figlio),
\textbf{gestione del flusso di informazioni} (ogni componente conosce solo gli aspetti del
data model che influenzano la sua view; il genitore ``spartisce'' i dati fra i figli).
Ma l'effetto centrale in React \`e:
\begin{importante}
\textbf{Se il valore di una prop cambia, React rigenera la view del componente}, invocando
nuovamente la funzione con i nuovi valori.
\end{importante}

\textbf{Parentesi: object destructuring.} Costrutto sintattico frequente nei componenti:
\texttt{const \{x, y\} = pippo} dichiara due variabili con i valori delle propriet\`a
omonime di \texttt{pippo}. Nei parametri formali: \texttt{function fun(\{x, y\}) \{\dots\}}
-- chi invoca passa un unico oggetto con propriet\`a \texttt{x} e \texttt{y}, disponibili
come variabili locali. Si applica anche agli array.

\section{Lo state: lo stato di un componente}

Con le sole props la view \`e ``statica'': i componenti cambiano al variare delle props, ma
non c'\`e modo di modificarle a partire dagli \emph{eventi UI}. Manca uno \textbf{stato} del
componente che cambi nel tempo e influenzi la view. Come conciliarlo con componenti che sono
funzioni \emph{pure}? (Se un listener modificasse una variabile locale -- possibile per via
della closure -- sarebbe un'operazione fine a s\'e stessa: al re-render la variabile verrebbe
ridichiarata e resettata.)
\begin{itemize}
  \item i componenti \emph{dichiarano} di aver bisogno di variabili di stato (tipo e valore
        iniziale);
  \item il valore corrente \`e conservato in un deposito \textbf{esterno ai componenti},
        gestito ``da React'';
  \item a ogni esecuzione, il componente ottiene da React il valore corrente, da usare nel
        determinare la view come gli altri input;
  \item ottiene anche una \textbf{funzione per modificare lo stato}, che \emph{non pu\`o
        essere invocata nel codice del componente} ma solo nelle callback asincrone: la
        modifica dello stato provoca il re-rendering (nuova esecuzione della funzione), e
        modificare lo stato dentro il componente causerebbe un \textbf{loop infinito}.
\end{itemize}
\begin{importante}
Come per le props: \textbf{se il valore di uno state cambia, React rigenera la view},
invocando di nuovo la funzione, che opera sul valore aggiornato.
\end{importante}

\subsection{Lo hook useState}
\begin{lstlisting}[language=JavaScript]
const [myState, setMyState] = useState<T>(valoreIniziale);
\end{lstlisting}
\begin{itemize}
  \item \texttt{useState<T>} prende il \emph{valore iniziale} (di tipo T), usato solo la
        prima volta (quando la variabile di stato non esiste ancora);
  \item restituisce una coppia (tecnicamente un array di due elementi -- la riga sopra \`e un
        \emph{array destructuring}): il \textbf{valore corrente} e la \textbf{funzione
        setter};
  \item il setter riceve un valore di tipo T che diventa il prossimo valore dello stato;
        pu\`o anche ricevere una \emph{funzione di aggiornamento} \texttt{(T) => T}, che
        determina il nuovo valore in base al precedente.
\end{itemize}
\begin{esempio}
\begin{lstlisting}[language=JavaScript]
type ContentNavItem = "negozio" | "libreria" | "lend" | "borrow" | "amici";
export function ContentNav(): ReactElement {
  const [activeItem, setActiveItem] = useState<ContentNavItem>("libreria");
  return (<nav className="content-nav">
    <div className={(activeItem === "negozio" ? ... )}
         onClick={() => setActiveItem("negozio")}>
      Negozio
    </div> ... </nav>);
}
\end{lstlisting}
\end{esempio}

\subsection{I React Hooks}
\texttt{useState} \`e uno \textbf{hook}: le funzioni della libreria React per ``agganciarsi''
al processo con cui React gestisce i componenti e influirvi parzialmente (il setter di
useState, oltre a modificare lo stato, \emph{forza il refresh}). Altri hook:
\texttt{useReducer}, \texttt{useContext}, \texttt{useRef}, \texttt{useMemo},
\texttt{useCallback}, \texttt{useEffect}\dots\ Riconoscibili dal nome che inizia per
\texttt{use}; si possono definire hook \emph{custom} a partire da quelli esistenti.
\begin{attenzione}
Gli hook \textbf{non} si possono usare ovunque: soltanto nello \emph{scope principale} di una
funzione-componente o di un custom hook.
\end{attenzione}

\section{Delega di un evento al genitore}

Talvolta lo stato di un componente deve aggiornarsi in seguito a eventi \emph{nei suoi
figli}. Esempio: una lista con selezione singola -- ogni voce \`e un componente che si
accorge del proprio click, ma per \emph{deselezionare le altre} serve la visione d'insieme,
posseduta solo dal componente-lista: sar\`a lui ad avere la variabile di stato che registra
la voce selezionata, e ogni figlio dovr\`a \emph{avvisarlo}.

\begin{definizione}[Delega al genitore]
Al verificarsi di un evento nel figlio, viene chiamata una \textbf{callback che appartiene al
genitore}. Il figlio dichiara di ricevere \emph{fra le sue props} una funzione definita dal
genitore: la callback da chiamare al verificarsi dell'evento. Per convenzione, le props di
gestione eventi hanno nome \texttt{on} + nome dell'evento. La delega pu\`o essere totale o
parziale (si pu\`o chiamare anche una callback del figlio, o una callback del figlio che
invoca quella del genitore): il figlio sceglie se ha cambiamenti di stato propri da attuare.
\end{definizione}

\textbf{Livelli di astrazione.} Quando un componente delega un evento si dice che lo
\emph{espone} al genitore; il nome dell'evento esposto va usato nel nome della prop
(evento ``action'' $\rightarrow$ prop \texttt{onAction}). Il figlio pu\`o esporre
\emph{direttamente} l'evento UI (la voce chiama la callback a ogni click: prop
\texttt{onClick} -- e in effetti anche gli ``elementi HTML'' in JSX sono componenti interni
di React, e i loro attributi sono props!), oppure \emph{elaborare} gli eventi UI e
determinare un evento a un \textbf{livello di astrazione pi\`u alto}: il componente-lista non
inoltra tutti i click (con posizione del mouse e tasti premuti), ma chiama la callback solo
quando \emph{cambia effettivamente la selezione} (tre click sulla stessa voce = una sola
selezione). Espone cos\`i un evento pi\`u alto, es. ``selection'' $\rightarrow$
\texttt{onSelection}: sono gli \textbf{eventi custom}, specifici dei componenti o
dell'applicazione.

\section{Form in React: binding bidirezionale}

In una SPA \emph{non} vogliamo il meccanismo automatico del browser per i form (provoca un
reload): li gestiamo con gli strumenti di React. \`E ``form-like'' qualunque componente
contenga elementi di input (INPUT, SELECT, TEXTAREA); anche i pulsanti di
conferma/annullamento possono essere qualsiasi elemento cliccabile. Passi:
\begin{enumerate}
  \item per ciascun elemento di input si dichiara uno \textbf{state} che ne contiene il
        valore: \texttt{<input type="text">} $\rightarrow$ string;
        \texttt{<select>} $\rightarrow$ il \texttt{value} della OPTION selezionata (string);
        \texttt{<input type="checkbox">} $\rightarrow$ boolean;
  \item si assegna all'input l'attributo \textbf{value}, facendolo corrispondere allo
        state: \texttt{<input value=\{stateVar\}>};
  \item si riceve l'evento \textbf{onChange} e lo si usa per modificare lo state:
\begin{lstlisting}[language=JavaScript]
<input type="text" value={stateVar}
       onChange={(e) => setStateVar(e.target.value)}>
\end{lstlisting}
\end{enumerate}
\`E insieme un property binding (stateVar passata per interpolazione) e un event binding
(stateVar settata in conseguenza di un evento): si parla di \textbf{doppio binding} o
\textbf{binding bidirezionale}.

\subsection{Tre tipologie di componenti form-like}
\begin{itemize}
  \item \textbf{Pannello di controllo}: sempre presente nella view; le impostazioni dei
        controlli si riflettono sulle view degli altri componenti. La delega al genitore
        avviene \emph{quasi a ogni cambiamento} di stato interno (che si riflette in uno
        stato del genitore, e quindi nelle props degli altri figli). Es.: filtri per
        categoria.
  \item \textbf{Modulo}: presente nella view; l'``azione'' si compie solo alla pressione del
        pulsante di conferma -- solo allora avviene la delega; gli altri eventi (incluso un
        eventuale reset) modificano solo lo stato interno. Es.: iscrizione a una newsletter
        in una colonna laterale.
  \item \textbf{Modal dialog}: \emph{non} presente nella view; appare per chiedere dati e
        resta finch\'e l'utente non conferma o annulla. La visibilit\`a \`e
        \emph{responsabilit\`a del genitore} (un componente invisibile non pu\`o cambiare il
        proprio stato); la delega avviene sia alla conferma sia all'annullamento, cos\`i il
        genitore pu\`o far scomparire la dialog. Es.: dialog per i dati di pagamento.
\end{itemize}

\section{Il ciclo di rendering di React}

\textbf{1. Inizializzazione.} Rendering di \texttt{App()}, quindi dei componenti: per
ciascuno React determina gli \emph{state}, le \emph{dipendenze} fra state+props del
componente e le props dei figli, ed effettua il rendering effettivo (incluso quello dei
figli).

\textbf{2. Cambiamenti di stato.} All'invocazione di \texttt{setVar} per una variabile di
stato \texttt{var}: React determina \emph{a quale componente} appartiene lo stato;
re-rendering del componente; React determina le conseguenze del cambiamento su altre
variabili/espressioni della funzione di rendering e \emph{quali props dei figli risultano
modificate}; re-rendering effettivo \textbf{solo di ci\`o che \`e cambiato}: elementi
montati/smontati, elementi con attributi cambiati, figli con props cambiate.

\textbf{Cosa provoca i cambiamenti di stato?} Le \emph{callback UI} (interazione con la UI) e
le \emph{funzioni async} (es. una fetch)\dots\ ovvero gli \textbf{effect}.

\section{useEffect: gli effetti collaterali del rendering}

\begin{definizione}[Effetto]
Un \textbf{effetto} \`e un blocco di codice che deve essere eseguito in conseguenza di certi
cambiamenti di props o state, ma che \emph{non pu\`o} essere eseguito direttamente nella
funzione-componente -- perch\'e influisce sul mondo esterno (violerebbe la purezza), perch\'e
pu\`o causare un re-rendering (es. cambia lo stato) o comunque interferire col rendering.
Viene eseguito come \emph{funzione asincrona} in seguito all'invocazione della funzione di
rendering. Esempi: \textbf{fetch di un dato necessario alla visualizzazione}; sottoscrizione
di un servizio esterno con callback; logging su DB locale; effetti di animazione.
\end{definizione}

\begin{lstlisting}[language=JavaScript]
useEffect(funzioneEffetto, dipendenzeEffetto?)
\end{lstlisting}
\begin{itemize}
  \item \texttt{funzioneEffetto}: la funzione che attua l'effetto;
  \item \texttt{dipendenzeEffetto}: array \emph{opzionale} delle variabili dello scope locale
        da cui l'effetto dipende (il cui cambiamento ne richiede la ri-esecuzione).
        \textbf{Array vuoto} $\Rightarrow$ eseguito solo al primo rendering;
        \textbf{array non specificato} $\Rightarrow$ eseguito \emph{a ogni} rendering.
\end{itemize}
La funzioneEffetto viene inserita nella \emph{coda degli eventi asincroni}: eseguita quando
il rendering \`e ormai terminato.

\subsection{useEffect + setState}
Essendo asincrona, la funzioneEffetto \emph{pu\`o} invocare un setter di stato -- attenzione
per\`o a non innescare un \textbf{loop infinito} fra cambi di stato e dipendenze:
\begin{itemize}
  \item caso semplice: il setState tocca una variabile che \emph{non} influisce sulle
        dipendenze dell'effetto $\Rightarrow$ nessuna ri-esecuzione;
  \item caso gestibile: la modifica influisce sulle dipendenze e l'effetto viene rieseguito
        -- allora la modifica non deve avvenire sistematicamente, ma sotto una condizione che
        prima o poi diventa falsa (un loop, ma finito, come un while).
\end{itemize}
Il caso pi\`u comune (spesso inevitabile): al primo rendering (e talvolta al cambio di una
prop) il componente deve fare una \textbf{fetch}; la fetch sta in un effetto F; all'arrivo
dei dati (then) li comunica con \texttt{setState(dato)}; segue il re-rendering. Bisogna
specificare le dipendenze di F in modo che il cambiamento di \texttt{dato} \emph{non ne
provochi la riesecuzione}: fra le dipendenze non ci deve essere \texttt{dato}, \textbf{n\'e
alcuna variabile locale calcolata a partire da esso}.

\subsection{La funzione di cleanup}
Alcuni effetti devono ``fare le pulizie'' prima di essere rieseguiti (es. terminare una
sottoscrizione o una connessione prima di aprirne un'altra, per evitare notifiche doppie o
memory leak). Inoltre un effetto pu\`o essere rieseguito per un re-rendering \emph{prima}
che l'esecuzione precedente si sia compiuta: es. fetch dipendente da una prop, prop
aggiornata, nuova fetch avviata prima che la precedente sia completata -- due fetch in corso
per due versioni del dato, ordine di arrivo imprevedibile: \emph{il dato obsoleto potrebbe
arrivare dopo e sovrascrivere quello corretto}. Infine, un componente pu\`o essere smontato
prima che un suo effetto si compia: l'effetto cercherebbe di modificare uno state che non
esiste pi\`u (React rimuove gli state dei componenti non pi\`u validi).

\begin{importante}[Cleanup]
La risposta \`e la \textbf{funzione di cleanup}: la funzioneEffetto deve \emph{restituire}
un'altra funzione (senza argomenti n\'e ritorno) che attua annullamento, invalidamento o
terminazione dell'effetto. React la invoca: \textbf{subito prima di reinvocare il medesimo
effetto} e \textbf{subito prima di smontare il componente}. \`E cos\`i importante che, in
fase di sviluppo, React invoca tutti gli effetti \emph{due volte di fila}: eventuali mancate
pulizie emergono pi\`u facilmente (anche se non sempre).
\end{importante}
